\SourceLeaf{23}
\[
\alpha\delta\xi=\delta\alpha.\int\!\int\frac{\mathrm{d}p.\mathrm{d}q.a\sin p}{r^3}\bigl(\frac{(a-x)x}{\alpha\alpha}+\frac{(b-y)y}{\beta\beta}+\frac{(c-z)z}{\gamma\gamma}\bigr)
\]
Nach Gleichung $[3]$ ist die rechte Seite dieser Gleichung entweder $=0$ oder $=-\frac{4\pi a\delta\alpha}{\alpha\beta\gamma}$, je nachdem, ob $M$ außerhalb oder innerhalb des Körpers liegt. Im ersten Fall gilt also
\[
[4]\qquad\delta\xi=0
\]
und im zweiten Fall
\[
[5]\qquad\delta\xi=-\frac{4\pi a\delta\alpha}{\alpha\alpha\beta\gamma}
\]
\par
Gleichung $[4]$ zeigt sofort, dass $\xi$ konstant ist oder dass die Anziehung $X$ der Masse proportional ist, und zwar für alle Ellipsoide, bei denen $\alpha\alpha-\beta\beta$ und $\alpha\alpha-\gamma\gamma$ konstant sind, deren drei Hauptschnitte also Ellipsen mit denselben Brennpunkten sind, solange der angezogene Punkt außerhalb des Sphäroids liegt. Da dieser Schluss auch dann streng gültig ist, wenn die Oberfläche des Sphäroids dem angezogenen Punkt beliebig nahe kommt, muss er sich auch auf das Sphäroid selbst erstrecken lassen, dessen Oberfläche durch den angezogenen Punkt geht.

Die Aufgabe, die Anziehung eines Sphäroids auf einen beliebigen äußeren Punkt zu bestimmen, lässt sich also auf zwei andere Aufgaben zurückführen. Zuerst sind die Abmessungen eines anderen Sphäroids zu bestimmen, das dieselben Brennpunkte wie das gegebene hat und durch den angezogenen Punkt geht. Sodann ist die Anziehung eines Sphäroids auf einen Punkt seiner Oberfläche zu bestimmen. Die erste Aufgabe hängt von der Lösung einer kubischen Gleichung ab. Es lässt sich leicht zeigen, dass sie stets genau eine reelle Wurzel hat, und es scheint überflüssig, hier bei ihr zu verweilen. Um aber die zweite Aufgabe zu lösen, betrachten wir den zweiten Fall, in dem der angezogene Punkt innerhalb des Körpers liegt. Da $\beta\beta=\alpha\alpha+BB-AA$, $\gamma\gamma=\alpha\alpha+CC-AA$ gilt, setzen wir diese Werte in Gleichung $5$ ein und setzen zugleich $\frac{A}{\alpha}=t$. Daraus ergibt sich
\[
\delta\xi=\frac{4a\pi tt\delta t}{A^3\BareRoot{((1-(1-\frac{BB}{AA})tt)(1-(1-\frac{CC}{AA})tt)}}
\]
oder, wenn wir wieder das Differentialzeichen $\mathrm{d}$ verwenden und integrieren,
\[
\xi=\frac{4a\pi}{A^3}\int\frac{tt\mathrm{d}t}{\BareRoot{((1-(1-\frac{BB}{AA})tt)(1-(1-\frac{CC}{AA})tt)}}
\]
Das Integral ist so zu nehmen, dass es für $t=0$ verschwindet, und dann für das bestimmte Sphäroid mit den Halbachsen $A,B,C$ bis $t=1$ zu erstrecken. Wir erhalten also
